\chapter{火山・マグマ}
\section{火砕密度流を考える背景}
\label{theory:volcanology:pyroflow}
噴火で破砕された粒子と高温のガスが混合し、周囲より重い混合物として斜面を流下するのが火砕密度流である。大気を取り込み、温度と濃度を変えながら進む。濃い流れと希薄なサージでは地形に対する挙動も異なる。堆積物は一般に分級が悪い塊状となりやすいが、全てが一様という意味ではない。

\section{ドーム崩落と噴煙柱からの発生}
\label{theory:volcanology:source}
溶岩ドームの崩落は岩塊と灰からなる流れを生む。メラピ1994年の論文では、谷を埋める流れと、それから分離して地形を横切る希薄サージが記録される。爆発的破砕と重力による崩落を分ける。噴煙柱の崩壊や大規模な噴出は別の発生経路であり、ドーム起源の小規模流とカルデラ形成を伴う大規模な活動を同一視しない。過去問のプレー型・メラピ型という歴史的呼称と大規模流の一般化は、個別解説で確認範囲を残す。\cite{volc_usgs_merapi}

\section{堆積分布は風と地形を分けて読む}
\label{theory:volcanology:distribution}
降下火山灰は風によって火口の風下へ運ばれ、細粒ほど遠くまで届きやすい。地面を流れる火砕流では谷、尾根、斜面変化、流れの濃度が経路を変えるため、舌状・ローブ状・扇状の分布や複数谷の埋積を考える。火口を一端とする単純な長円形を一般的な火砕流の分布として決め付けない。元問題図の転載をせず、移動媒体と支配条件の違いを読み取る。\cite{volc_usgs_msh_flow,volc_usgs_merapi}

\section{固体の内部で局所的に融解する}\label{theory:volcano:melting}
マントルの多くは固体である。地震波を伝える固体であることと、地質学的な時間にゆっくり流動することは両立する。マグマは地球内部を一面に覆う液体ではなく、温度・圧力・組成の経路が融解条件を越える場所で生じる。
多成分岩石では融解開始の固相線と全てが融ける液相線が異なる。間では固相とメルトが共存し、部分溶融となる。初生マグマは原岩の組成をそのまま全量コピーしたものではない\cite{vol50_melt_ref}。
物理での相図の読み方を使うと、融解開始は概念的に \(T_{\rm path}(P)>T_{\rm solidus}(P,X)\) と表せる。上昇するマントルでは圧力が下がり、ほぼ断熱的な温度経路と固相線の交差で減圧融解が始まる。海嶺の成因を「低温の物質が突然加熱される」と説明しない。ホットスポットの熱異常と上昇減圧は別の要因である。
沈み込みでは含水物質の脱水により流体がマントルウェッジへ供給され、融解開始温度を下げる\cite{vol50_subduction_ref}。温度を上げる、圧力を下げる、水等で固相線を下げるという三経路を分ける。
メルトが周囲より低密度なら浮力で上昇し得る。概略の駆動力は体積 \(V\) 当たり \((\rho_{\rm rock}-\rho_{\rm melt})gV\) だが、実際の移動には割れ目・応力・粘性・冷却も関係する。全マグマだまりを中立浮力だけで決定しない。

\section{原岩と別マグマのどちらと比較するか}
初生的なマントル融解液は代表的に玄武岩質である。結晶が分離すれば残る液体の組成は変わり、地殻の融解や同化・混合も別のマグマを作る\cite{vol50_melt_ref}。
玄武岩質は流紋岩質等よりSiO2に乏しい側だが、原岩マントルに比べればNa・K等の不適合元素がメルトへ入り濃集する場合がある。「富む」は比較対象を明記する。
JMAの1999年表では玄武岩質の代表噴出温度が1000--1200℃、流紋岩質が700--900℃で、珪酸分が多い側は一般に粘性が大きい\cite{jma_magma1999}。これは原典時点の分類と代表値であり、全組成・全含水量・全結晶量の固定境界ではない。

\section{海溝・前弧・火山弧・背弧}\label{theory:volcano:arc}
地表で海溝から陸側へ前弧・火山弧・背弧を並べ、地下では沈み込むスラブとその上のマントルウェッジを対応させる。火山前線は火山の分布の海溝側限界で、海溝そのものと違う\cite{usgs_plate_boundaries}。
火山が少ない前弧でもプレート境界やスラブの地震活動があり、「火山がないから地震もない」は誤りである。背弧に拡大海盆ができる場合はあるが、全島弧が同じ拡大状態を持つとはしない。
アルカリ元素と水を一つの変数として扱わない。島弧の背弧方向のアルカリ濃度には一般的な増加傾向を学ぶ一方、GSJの東北日本弧の発表には含水量が背弧へ減少する分布の報告がある\cite{vol50_arcwater_ref}。融解度・供給源・脱ガスの違いを考え、全弧の水濃度へ固定した順位を使わない。
中央海嶺は発散境界で新しい海洋地殻を作る。USGS概説では毎年の地球の噴出物の約4分の3が拡大海嶺起源と推計される\cite{vol50_ridge_ref}。陸上噴火の個数や被害の印象と、全世界の体積率を区別する。

\section{間欠爆発と連続溶岩噴泉}\label{theory:volcano:styles}
低粘性の玄武岩質マグマの連続的な噴泉や流出はハワイ式の代表である。ストロンボリ式はガスの断続的な解放により小爆発を繰り返し、火山弾・スコリアを放出する\cite{jma_magma1999,vol50_styles_ref}。同じ火山が常に一つの様式だけで噴火するわけではない。
高粘性の溶岩は流れにくく火口近くへ蓄積してドームを作り得る。ガスが抜ける時間、供給量、結晶量等も噴火様式へ関係する。岩種名だけで爆発の有無を必ず決めない。
マグマ水蒸気爆発では外来水とマグマの急な熱交換・蒸気膨張・破砕が関与する。水蒸気噴火で既存岩石が破砕される場合と、新しいマグマが外来水へ接触する場合を区別する。

\section{溶岩表面の形と冷却}\label{theory:volcano:lava}
パホイホイは代表的に玄武岩質の流動的な溶岩で、滑らか・縄状等の表面を持つ。アアは粗いクリンカーに覆われる。同じ玄武岩でも冷却・結晶化・変形速度等で一方から他方へ変化し得る\cite{vol50_lava_ref}。滑らかだから流紋岩という対応は誤りである。
ブロック溶岩では大きい岩塊の比較的平滑な破断面が特徴となる。表面を形態で分ける軸と、岩石の化学組成を分ける軸を混同しない。
厚さ／長さ等のアスペクト比は地形を比較する指標である。原典の1/8区分は分類上の目安として保存し、普遍的な相転移や流動停止の定数とはしない。長さの測り方・流量・地形・冷却の影響を考える。

\section{ガスの主要成分と平衡条件}\label{theory:volcano:gases}
通常の火山ガスの最多成分は水蒸気であり、CO2、SO2、H2S、水素ハロゲン化物等も放出される\cite{vol50_gas_ref}。危険性の大きさと濃度順位は別の情報である。
減圧ではメルト中の揮発性分の溶解度が下がり、気泡が分離し膨張する。CO2と水の脱ガス段階は同じではない\cite{vol50_degassing_ref}。閉塞で過圧が生じる場合と、開いた系でガスが抜ける場合を分ける。外圧低下だけからガス絶対圧力の単調増加を推論しない。
高温マグマのガス、冷却された噴気、地下水と反応した熱水ガスは組成が違い得る。CH4を通常高温の主要ガスの比較で選外とすることは、全自然条件で存在しないことを意味しない。令和7年度の100気圧・1200℃・玄武岩質という条件も、酸化還元状態・全元素量なしに任意の正確なモル割合を作れるわけではない。

\section{粒子とガス、水、岩屑を分ける}\label{theory:volcano:hazards}
火砕流は高温のガス・粒子混合、火山泥流は水・粒子混合、岩屑なだれは崩壊岩体の移動を中心に考える。速さの目安と材料を同時に確認する。噴煙柱の崩壊やドーム崩落と、融雪・湖水・降雨による泥流形成を同一視しない。火砕密度流の背景は第\ref{theory:volcanology:pyroflow}節へ戻る。
大量の物質が海へ入り水を急激に動かすと津波が生じる場合がある。噴火・海底地殻変動・崩壊による水塊の変位を考え、全噴火が津波を起こすとはしない。
噴煙が高くても、灰の短期的な遮蔽とSO2由来硫酸エアロゾルの長距離・長期間の気候影響を区別する\cite{vol50_climate_ref}。粒子の沈降時間、化学反応、成層圏での滞留という物理と化学の段階を追う。

\section{火炎状構造の「火」は成因ではない}\label{theory:volcano:terms}
カルデラ、火山岩の安山岩、溶岩片が付着したアグルチネート、マグマ水蒸気爆発は火山の直接用語である。
火炎状構造は未固結の軽い堆積層が上の重い層の間へ押し上げられる等の変形を指し、名称は形の比喩である\cite{vol50_terms_ref}。堆積物の形成地域が火山である可能性を認めても、変形の直接機構を炎や噴火としない。名前を覚えた後に、物質・応力・流れの機構を説明できることを学習目標とする。


\section{過去問と個別の解法}
\section{島弧沈み込み系と中央海嶺の用語を区別する。}

\subsection{平成25年度 問26}
\label{q:h25-reg-applied-26}

\noindent\textbf{原典：} sources/H25/17\_applied\_science.pdf、PDF物理10頁。

\textbf{公式正答：} 5（single）。

\paragraph{問題の要旨} 島弧沈み込み系と中央海嶺の用語を区別する。

\Needspace{4\baselineskip}\paragraph{確認した条件・記述（独自要約）}\mbox{}\par

\noindent\textbf{原典確認：} S7\_050の29原典画像で34問の全五肢・否定・数値を照合。\par

\noindent\textbf{対象：} 島弧沈み込み系と中央海嶺の用語を区別する。\par

\paragraph{問われる核心} 島弧沈み込み系と中央海嶺の用語を区別する。

\paragraph{必要知識} 海溝から陸側へ前弧・火山弧・背弧を並べ、スラブとその上のマントルウェッジを考える。中央海嶺は発散境界であり島弧沈み込み系と直接の用語が違うので5。背弧海盆にも拡大軸があり得るが、地球規模の中央海嶺を島弧そのものへ交換しない。

\paragraph{考え方・正答を選ぶ理由（査読待ち初稿）} 海溝から陸側へ前弧・火山弧・背弧を並べ、スラブとその上のマントルウェッジを考える。中央海嶺は発散境界であり島弧沈み込み系と直接の用語が違うので5。背弧海盆にも拡大軸があり得るが、地球規模の中央海嶺を島弧そのものへ交換しない。

\Needspace{4\baselineskip}\paragraph{選択肢の検討（査読待ち初稿）}\mbox{}\par

\noindent\textbf{1：代表的条件で適切} マントルウェッジ。沈み込むスラブの上のくさび状マントルで、流体供給と島弧マグマ生成に関係する。\par

\noindent\textbf{2：代表的条件で適切} 背弧。火山弧の海溝と反対側で、伸張する背弧盆地等が形成され得る。\par

\noindent\textbf{3：代表的条件で適切} 火山前線。島弧火山の分布の海溝側限界を示す線。海溝そのものと違う。\par

\noindent\textbf{4：代表的条件で適切} 沈み込み帯。収束するプレートが地下へ沈み込む島弧活動の主要な生成場。\par

\noindent\textbf{5：不適切} 中央海嶺。プレートが離れ、新しい海洋地殻を作る発散境界。島弧の直接構成用語ではない。\par

\paragraph{再利用する考え方} 地質環境・生成機構・物質の名称を対応させ、平均値と全称、典型例と例外、当時の分類と現在の知見を分ける。

\paragraph{共通理論} \hyperref[theory:volcano:arc]{第\ref*{theory:volcano:arc}節}

\par 参考：\cite{usgs_plate_boundaries}。

\par 難易度 3（校準初稿）、優先度未評価。

\subsection{令和元年度再試験 問23}
\label{q:r01-ret-applied-23}

\noindent\textbf{原典：} sources/R01\_retest/17\_applied\_science.pdf、PDF物理8頁。

\textbf{公式正答：} 3（single）。

\par\textbf{検討未確定：} 公式3。背弧方向のアルカリと水の濃度は別の変数で、水まで一般に増加と改変しない。全マグマの上部マントル起源という全称と温度の全噴火への適用も留保。アルカリ傾向の原論文による追加照合が必要。

\paragraph{問題の要旨} 火山前線・背弧の組成傾向とマグマ温度を比較する。

\Needspace{4\baselineskip}\paragraph{確認した条件・記述（独自要約）}\mbox{}\par

\noindent\textbf{原典確認：} S7\_050の29原典画像で34問の全五肢・否定・数値を照合。\par

\noindent\textbf{対象：} 火山前線・背弧の組成傾向とマグマ温度を比較する。\par

\paragraph{問われる核心} 火山前線・背弧の組成傾向とマグマ温度を比較する。

\paragraph{必要知識} 火山前線の海溝と反対側へ進むと、玄武岩中のアルカリ濃度が一般に高まる傾向がある。原典3は小さくなるとするため不適切。ただし水の濃度を常にアルカリと同方向としない。GSJの東北日本弧の調査には含水量が背弧へ減る分布の報告があり、アルカリと水を分けて比較する。

\paragraph{考え方・正答を選ぶ理由（査読待ち初稿）} 火山前線の海溝と反対側へ進むと、玄武岩中のアルカリ濃度が一般に高まる傾向がある。原典3は小さくなるとするため不適切。ただし水の濃度を常にアルカリと同方向としない。GSJの東北日本弧の調査には含水量が背弧へ減る分布の報告があり、アルカリと水を分けて比較する。

\Needspace{4\baselineskip}\paragraph{選択肢の検討（査読待ち初稿）}\mbox{}\par

\noindent\textbf{1：条件・解釈を留保} マントル部分融解によるマグマ。初生的な玄武岩質マグマの主要な生成経路。地殻の融解からもマグマが作られるので全種類の唯一の経路とはしない。\par

\noindent\textbf{2：代表的条件で適切} 海嶺・ホットスポットは主に玄武岩質。代表的な初生噴出物は玄武岩質。分化した溶岩等の例外はある。\par

\noindent\textbf{3：不適切} 背弧へアルカリと水が系統的に低下。島弧の代表的アルカリ濃度は背弧へ増す傾向で、アルカリと水を同方向の普遍則へ結び付けた記述は適切でない。水の傾向には別の検証が必要。\par

\noindent\textbf{4：代表的条件で適切} 島弧にはシリカに富む分化マグマもある。安山岩・デイサイト・流紋岩等も噴出し、初生玄武岩だけではない。\par

\noindent\textbf{5：代表的条件で適切} 玄武岩1000–1200℃・流紋岩700–900℃。JMAの1999年表にも記載された代表的な噴出温度範囲で、全噴火の厳密な境界ではない。\par

\paragraph{再利用する考え方} 地質環境・生成機構・物質の名称を対応させ、平均値と全称、典型例と例外、当時の分類と現在の知見を分ける。

\paragraph{共通理論} \hyperref[theory:volcano:arc]{第\ref*{theory:volcano:arc}節}

\par 参考：\cite{vol50_melt_ref,jma_magma1999,vol50_arcwater_ref}。

\par 難易度 4（校準初稿）、優先度未評価。

\subsection{令和6年度 問25}
\label{q:r06-reg-applied-25}

\noindent\textbf{原典：} sources/R06/17\_applied\_science.pdf、PDF物理8頁。

\textbf{公式正答：} 5（single）。

\par\textbf{検討未確定：} 原典4の島弧・陸弧を海の有無だけで定義する単純化と、前弧に火山が全くないという全称は留保。公式5の地震活動の誤りを保持。

\paragraph{問題の要旨} 海溝・前弧・火山弧・背弧と地震活動を整理する。

\Needspace{4\baselineskip}\paragraph{確認した条件・記述（独自要約）}\mbox{}\par

\noindent\textbf{原典確認：} S7\_050の29原典画像で34問の全五肢・否定・数値を照合。\par

\noindent\textbf{対象：} 海溝・前弧・火山弧・背弧と地震活動を整理する。\par

\paragraph{問われる核心} 海溝・前弧・火山弧・背弧と地震活動を整理する。

\paragraph{必要知識} 沈み込むプレートを地下まで追うと、火山前線より海溝側でも地震が起きる。火山分布と地震分布の範囲を同一視した5が不適切。海溝→前弧→火山弧→背弧の配列を地表と断面で対応させる。

\paragraph{考え方・正答を選ぶ理由（査読待ち初稿）} 沈み込むプレートを地下まで追うと、火山前線より海溝側でも地震が起きる。火山分布と地震分布の範囲を同一視した5が不適切。海溝→前弧→火山弧→背弧の配列を地表と断面で対応させる。

\Needspace{4\baselineskip}\paragraph{選択肢の検討（査読待ち初稿）}\mbox{}\par

\noindent\textbf{1：代表的条件で適切} 海洋側に海溝。島弧海溝の基本配列では海溝が沈み込みの海洋側にある。\par

\noindent\textbf{2：代表的条件で適切} 陸側へ前弧・火山弧・背弧。海溝から火山前線までの前弧、火山弧、その背後の背弧を順に読む。\par

\noindent\textbf{3：代表的条件で適切} 背弧で海が形成され得る。伸張して背弧盆地・海が形成される場合がある。必ず拡大海盆を持つわけではない。\par

\noindent\textbf{4：条件・解釈を留保} 背後の海による島弧・陸弧の呼称。過去問の基本図式は大陸から海で隔たる島弧と大陸縁の陸弧。海の有無だけで全遷移形態を排他的に分類する厳密定義には留保する。\par

\noindent\textbf{5：不適切} 前弧は地震活動も静穏。沈み込み境界やスラブに関係する地震があり、火山が少ないことから地震も静穏とは言えない。\par

\paragraph{再利用する考え方} 地質環境・生成機構・物質の名称を対応させ、平均値と全称、典型例と例外、当時の分類と現在の知見を分ける。

\paragraph{共通理論} \hyperref[theory:volcano:arc]{第\ref*{theory:volcano:arc}節}

\par 参考：\cite{usgs_plate_boundaries,vol50_subduction_ref}。

\par 難易度 4（校準初稿）、優先度未評価。

\section{ハワイ式とストロンボリ式・火砕流を区別する。}

\subsection{平成25年度 問25}
\label{q:h25-reg-applied-25}

\noindent\textbf{原典：} sources/H25/17\_applied\_science.pdf、PDF物理9頁。

\textbf{公式正答：} 3（single）。

\paragraph{問題の要旨} ハワイ式とストロンボリ式・火砕流を区別する。

\Needspace{4\baselineskip}\paragraph{確認した条件・記述（独自要約）}\mbox{}\par

\noindent\textbf{原典確認：} S7\_050の29原典画像で34問の全五肢・否定・数値を照合。\par

\noindent\textbf{対象：} ハワイ式とストロンボリ式・火砕流を区別する。\par

\paragraph{問われる核心} ハワイ式とストロンボリ式・火砕流を区別する。

\paragraph{必要知識} マグマの粘性とガス放出の仕方を対応させる。ハワイ式の代表は低粘性の玄武岩質で、流紋岩とした3が誤り。ストロンボリ式の間欠爆発、火砕流の高速なガス粒子混合、外来水との爆発、ドーム形成をそれぞれ別の現象として整理する。

\paragraph{考え方・正答を選ぶ理由（査読待ち初稿）} マグマの粘性とガス放出の仕方を対応させる。ハワイ式の代表は低粘性の玄武岩質で、流紋岩とした3が誤り。ストロンボリ式の間欠爆発、火砕流の高速なガス粒子混合、外来水との爆発、ドーム形成をそれぞれ別の現象として整理する。

\Needspace{4\baselineskip}\paragraph{選択肢の検討（査読待ち初稿）}\mbox{}\par

\noindent\textbf{1：代表的条件で適切} 間欠的な火山弾・スコリア放出。ストロンボリ式は比較的低粘性のマグマからガスが断続的に解放される小爆発を伴う。\par

\noindent\textbf{2：代表的条件で適切} 高温ガスと火砕物の高速流下。火砕流は高温のガス・粒子混合流で、地表を高速に移動する。\par

\noindent\textbf{3：不適切} 低粘性流紋岩のハワイ式。典型的なハワイ式は流動性に富む玄武岩質マグマの溶岩噴泉等。流紋岩へ交換した点が誤り。\par

\noindent\textbf{4：代表的条件で適切} 外来水とマグマの爆発的相互作用。マグマ水蒸気爆発では急な熱交換・水の蒸発・破砕が関与する。マグマ由来の揮発性分だけの発泡と分ける。\par

\noindent\textbf{5：代表的条件で適切} 高粘性溶岩のドーム。安山岩質～デイサイト質等の高粘性溶岩は火口近くに堆積してドームを作り得る。組成だけで全ての形を固定しない。\par

\paragraph{再利用する考え方} 地質環境・生成機構・物質の名称を対応させ、平均値と全称、典型例と例外、当時の分類と現在の知見を分ける。

\paragraph{共通理論} \hyperref[theory:volcano:styles]{第\ref*{theory:volcano:styles}節}

\par 参考：\cite{jma_magma1999,vol50_styles_ref}。

\paragraph{関連問題} \hyperref[q:h30-reg-applied-29]{H30-REG-APPLIED-29}

\par 難易度 3（校準初稿）、優先度未評価。

\subsection{平成30年度 問29}
\label{q:h30-reg-applied-29}

\noindent\textbf{原典：} sources/H30/17\_applied\_science.pdf、PDF物理10頁。

\textbf{公式正答：} 3（single）。

\paragraph{問題の要旨} ハワイ式とストロンボリ式・火砕流を区別する。

\Needspace{4\baselineskip}\paragraph{確認した条件・記述（独自要約）}\mbox{}\par

\noindent\textbf{原典確認：} S7\_050の29原典画像で34問の全五肢・否定・数値を照合。\par

\noindent\textbf{対象：} ハワイ式とストロンボリ式・火砕流を区別する。\par

\paragraph{問われる核心} ハワイ式とストロンボリ式・火砕流を区別する。

\paragraph{必要知識} マグマの粘性とガス放出の仕方を対応させる。ハワイ式の代表は低粘性の玄武岩質で、流紋岩とした3が誤り。ストロンボリ式の間欠爆発、火砕流の高速なガス粒子混合、外来水との爆発、ドーム形成をそれぞれ別の現象として整理する。

\paragraph{考え方・正答を選ぶ理由（査読待ち初稿）} マグマの粘性とガス放出の仕方を対応させる。ハワイ式の代表は低粘性の玄武岩質で、流紋岩とした3が誤り。ストロンボリ式の間欠爆発、火砕流の高速なガス粒子混合、外来水との爆発、ドーム形成をそれぞれ別の現象として整理する。

\Needspace{4\baselineskip}\paragraph{選択肢の検討（査読待ち初稿）}\mbox{}\par

\noindent\textbf{1：代表的条件で適切} 間欠的な火山弾・スコリア放出。ストロンボリ式は比較的低粘性のマグマからガスが断続的に解放される小爆発を伴う。\par

\noindent\textbf{2：代表的条件で適切} 高温ガスと火砕物の高速流下。火砕流は高温のガス・粒子混合流で、地表を高速に移動する。\par

\noindent\textbf{3：不適切} 低粘性流紋岩のハワイ式。典型的なハワイ式は流動性に富む玄武岩質マグマの溶岩噴泉等。流紋岩へ交換した点が誤り。\par

\noindent\textbf{4：代表的条件で適切} 外来水とマグマの爆発的相互作用。マグマ水蒸気爆発では急な熱交換・水の蒸発・破砕が関与する。マグマ由来の揮発性分だけの発泡と分ける。\par

\noindent\textbf{5：代表的条件で適切} 高粘性溶岩のドーム。安山岩質～デイサイト質等の高粘性溶岩は火口近くに堆積してドームを作り得る。組成だけで全ての形を固定しない。\par

\paragraph{再利用する考え方} 地質環境・生成機構・物質の名称を対応させ、平均値と全称、典型例と例外、当時の分類と現在の知見を分ける。

\paragraph{共通理論} \hyperref[theory:volcano:styles]{第\ref*{theory:volcano:styles}節}

\par 参考：\cite{jma_magma1999,vol50_styles_ref}。

\paragraph{関連問題} \hyperref[q:h25-reg-applied-25]{H25-REG-APPLIED-25}

\par 難易度 3（校準初稿）、優先度未評価。

\subsection{令和元年度再試験 問24}
\label{q:r01-ret-applied-24}

\noindent\textbf{原典：} sources/R01\_retest/17\_applied\_science.pdf、PDF物理8頁。

\textbf{公式正答：} 5（single）。

\paragraph{問題の要旨} 噴火現象を気体・水・溶岩・岩屑の移動へ分ける。

\Needspace{4\baselineskip}\paragraph{確認した条件・記述（独自要約）}\mbox{}\par

\noindent\textbf{原典確認：} S7\_050の29原典画像で34問の全五肢・否定・数値を照合。\par

\noindent\textbf{対象：} 噴火現象を気体・水・溶岩・岩屑の移動へ分ける。\par

\paragraph{問われる核心} 噴火現象を気体・水・溶岩・岩屑の移動へ分ける。

\paragraph{必要知識} 流れる物質と速度を確認する。火砕流は高温ガスと火砕粒子の混合が高速に流下する現象であり5の低速度が誤り。外来水を含む爆発・泥流、液状溶岩、崩壊する岩屑の移動はそれぞれ異なる。

\paragraph{考え方・正答を選ぶ理由（査読待ち初稿）} 流れる物質と速度を確認する。火砕流は高温ガスと火砕粒子の混合が高速に流下する現象であり5の低速度が誤り。外来水を含む爆発・泥流、液状溶岩、崩壊する岩屑の移動はそれぞれ異なる。

\Needspace{4\baselineskip}\paragraph{選択肢の検討（査読待ち初稿）}\mbox{}\par

\noindent\textbf{1：代表的条件で適切} 液状マグマの流出。液体として流れ出て移動する溶岩流。気体粒子混合の火砕流と分ける。\par

\noindent\textbf{2：代表的条件で適切} ストロンボリの間欠爆発。比較的低粘性のマグマから気泡が断続的に放出される。\par

\noindent\textbf{3：代表的条件で適切} 外来水の急な水蒸気化。マグマと水の急な熱交換で蒸気の膨張・破砕が生じ得る。\par

\noindent\textbf{4：代表的条件で適切} 崩壊山体の岩屑なだれ。不安定な火山体が崩れ高速に移動する場合がある。\par

\noindent\textbf{5：不適切} 火砕流は時速数km以下。一般的な火砕流は時速数十kmから100km以上等の高速で、ゆっくりした数km以下の流れという説明が不適切。\par

\paragraph{再利用する考え方} 地質環境・生成機構・物質の名称を対応させ、平均値と全称、典型例と例外、当時の分類と現在の知見を分ける。

\paragraph{共通理論} \hyperref[theory:volcano:styles]{第\ref*{theory:volcano:styles}節}

\par 参考：\cite{jma_magma1999,volc_usgs_msh_flow}。

\paragraph{関連問題} \hyperref[q:r03-reg-applied-26]{R03-REG-APPLIED-26}

\par 難易度 3（校準初稿）、優先度未評価。

\subsection{令和3年度 問26}
\label{q:r03-reg-applied-26}

\noindent\textbf{原典：} sources/R03/17\_applied\_science.pdf、PDF物理9頁。

\textbf{公式正答：} 3（single）。

\paragraph{問題の要旨} ストロンボリの間欠爆発と連続溶岩噴泉を区別する。

\Needspace{4\baselineskip}\paragraph{確認した条件・記述（独自要約）}\mbox{}\par

\noindent\textbf{原典確認：} S7\_050の29原典画像で34問の全五肢・否定・数値を照合。\par

\noindent\textbf{対象：} ストロンボリの間欠爆発と連続溶岩噴泉を区別する。\par

\paragraph{問われる核心} ストロンボリの間欠爆発と連続溶岩噴泉を区別する。

\paragraph{必要知識} 間欠か連続かを原典から読む。ストロンボリ式は気泡の解放による間欠的小爆発で、割れ目からの連続的噴泉とした3が誤り。平成・令和の似た五肢でも、誤りが岩種か速度か噴出様式かを確かめる。

\paragraph{考え方・正答を選ぶ理由（査読待ち初稿）} 間欠か連続かを原典から読む。ストロンボリ式は気泡の解放による間欠的小爆発で、割れ目からの連続的噴泉とした3が誤り。平成・令和の似た五肢でも、誤りが岩種か速度か噴出様式かを確かめる。

\Needspace{4\baselineskip}\paragraph{選択肢の検討（査読待ち初稿）}\mbox{}\par

\noindent\textbf{1：代表的条件で適切} 時速数十から100km以上の火砕流。高温のガス粒子混合流の代表的高速移動。速度の全現象一律境界ではない。\par

\noindent\textbf{2：代表的条件で適切} 液状溶岩流。破片の混合ではなく液体としてのマグマが流出する。\par

\noindent\textbf{3：不適切} 連続噴泉をストロンボリ式とする。連続的な溶岩噴泉は代表的にハワイ式。ストロンボリは間欠的な爆発が特徴。\par

\noindent\textbf{4：代表的条件で適切} 山体崩壊の岩屑なだれ。火山体の不安定部が崩れ高速に流下する。\par

\noindent\textbf{5：代表的条件で適切} 外来水の相変化を伴う爆発。マグマ水蒸気爆発で急な熱交換・膨張等が起こる。\par

\paragraph{再利用する考え方} 地質環境・生成機構・物質の名称を対応させ、平均値と全称、典型例と例外、当時の分類と現在の知見を分ける。

\paragraph{共通理論} \hyperref[theory:volcano:styles]{第\ref*{theory:volcano:styles}節}

\par 参考：\cite{jma_magma1999,vol50_styles_ref,volc_usgs_msh_flow}。

\paragraph{関連問題} \hyperref[q:r01-ret-applied-24]{R01-RET-APPLIED-24}

\par 難易度 3（校準初稿）、優先度未評価。

\section{火山ガスの主要成分と温度・酸化還元条件を区別する。}

\subsection{平成28年度 問28}
\label{q:h28-reg-applied-28}

\noindent\textbf{原典：} sources/H28/17\_applied\_science.pdf、PDF物理11頁。

\textbf{公式正答：} 5（single）。

\par\textbf{検討未確定：} 理論的平均と主要の厳密な濃度基準は未提示。特にR07の100気圧・1200℃平衡の定量表と酸化還元条件の照合が残る。通常の火山ガスの代表比較として説明し、CH4の全条件不在やH2の固定割合を主張しない。

\paragraph{問題の要旨} 火山ガスの主要成分と温度・酸化還元条件を区別する。

\Needspace{4\baselineskip}\paragraph{確認した条件・記述（独自要約）}\mbox{}\par

\noindent\textbf{原典確認：} S7\_050の29原典画像で34問の全五肢・否定・数値を照合。\par

\noindent\textbf{対象：} 火山ガスの主要成分と温度・酸化還元条件を区別する。\par

\paragraph{問われる核心} 火山ガスの主要成分と温度・酸化還元条件を区別する。

\paragraph{必要知識} 水蒸気、CO2、硫黄ガス、水素ハロゲン化物等が通常の火山ガスの代表で、候補中ではCH4を選ぶ。高温マグマのガスと低温熱水・堆積物由来のガスを分ける。令和7年度は100気圧・1200℃・玄武岩質の平衡条件が原典にあり、平成28年度・令和元再試験へその条件を遡って追加しない。

\paragraph{考え方・正答を選ぶ理由（査読待ち初稿）} 水蒸気、CO2、硫黄ガス、水素ハロゲン化物等が通常の火山ガスの代表で、候補中ではCH4を選ぶ。高温マグマのガスと低温熱水・堆積物由来のガスを分ける。令和7年度は100気圧・1200℃・玄武岩質の平衡条件が原典にあり、平成28年度・令和元再試験へその条件を遡って追加しない。

\Needspace{4\baselineskip}\paragraph{選択肢の検討（査読待ち初稿）}\mbox{}\par

\noindent\textbf{1：代表的条件で適切} H2O。通常の火山ガスの最多成分は水蒸気であり、酸性有毒ガスが量の大部分という意味ではない。\par

\noindent\textbf{2：代表的条件で適切} CO2。マグマから放出される代表的な炭素ガス。水蒸気との量比は環境に依存する。\par

\noindent\textbf{3：代表的条件で適切} SO2。硫黄ガスの代表で酸化的な高温火山ガスに重要。冷却・熱水系では化学種や吸収が変わる。\par

\noindent\textbf{4：代表的条件で適切} H2S。硫黄の還元的な化学種であり、火山ガス・熱水系に含まれる。\par

\noindent\textbf{5：不適切} CH4。通常の高温マグマに由来する平均主要ガスの比較では適さない。低温・還元環境の存在を含め、厳密に常に零とする意味ではない。\par

\paragraph{再利用する考え方} 地質環境・生成機構・物質の名称を対応させ、平均値と全称、典型例と例外、当時の分類と現在の知見を分ける。

\paragraph{共通理論} \hyperref[theory:volcano:gases]{第\ref*{theory:volcano:gases}節}

\par 参考：\cite{vol50_gas_ref}。

\paragraph{関連問題} \hyperref[q:r01-ret-applied-22]{R01-RET-APPLIED-22}

\par 難易度 4（校準初稿）、優先度未評価。

\subsection{令和元年度再試験 問22}
\label{q:r01-ret-applied-22}

\noindent\textbf{原典：} sources/R01\_retest/17\_applied\_science.pdf、PDF物理7頁。

\textbf{公式正答：} 4（single）。

\par\textbf{検討未確定：} 理論的平均と主要の厳密な濃度基準は未提示。特にR07の100気圧・1200℃平衡の定量表と酸化還元条件の照合が残る。通常の火山ガスの代表比較として説明し、CH4の全条件不在やH2の固定割合を主張しない。

\paragraph{問題の要旨} 火山ガスの主要成分と温度・酸化還元条件を区別する。

\Needspace{4\baselineskip}\paragraph{確認した条件・記述（独自要約）}\mbox{}\par

\noindent\textbf{原典確認：} S7\_050の29原典画像で34問の全五肢・否定・数値を照合。\par

\noindent\textbf{対象：} 火山ガスの主要成分と温度・酸化還元条件を区別する。\par

\paragraph{問われる核心} 火山ガスの主要成分と温度・酸化還元条件を区別する。

\paragraph{必要知識} 水蒸気、CO2、硫黄ガス、水素ハロゲン化物等が通常の火山ガスの代表で、候補中ではCH4を選ぶ。高温マグマのガスと低温熱水・堆積物由来のガスを分ける。令和7年度は100気圧・1200℃・玄武岩質の平衡条件が原典にあり、平成28年度・令和元再試験へその条件を遡って追加しない。

\paragraph{考え方・正答を選ぶ理由（査読待ち初稿）} 水蒸気、CO2、硫黄ガス、水素ハロゲン化物等が通常の火山ガスの代表で、候補中ではCH4を選ぶ。高温マグマのガスと低温熱水・堆積物由来のガスを分ける。令和7年度は100気圧・1200℃・玄武岩質の平衡条件が原典にあり、平成28年度・令和元再試験へその条件を遡って追加しない。

\Needspace{4\baselineskip}\paragraph{選択肢の検討（査読待ち初稿）}\mbox{}\par

\noindent\textbf{1：代表的条件で適切} CO2。マグマから放出される代表的な炭素ガス。水蒸気との量比は環境に依存する。\par

\noindent\textbf{2：代表的条件で適切} H2S。硫黄の還元的な化学種であり、火山ガス・熱水系に含まれる。\par

\noindent\textbf{3：代表的条件で適切} HCl。火山ガス中の代表的な水素ハロゲン化物。地下水への吸収で地表の濃度は変わり得る。\par

\noindent\textbf{4：不適切} CH4。通常の高温マグマに由来する平均主要ガスの比較では適さない。低温・還元環境の存在を含め、厳密に常に零とする意味ではない。\par

\noindent\textbf{5：代表的条件で適切} SO2。硫黄ガスの代表で酸化的な高温火山ガスに重要。冷却・熱水系では化学種や吸収が変わる。\par

\paragraph{再利用する考え方} 地質環境・生成機構・物質の名称を対応させ、平均値と全称、典型例と例外、当時の分類と現在の知見を分ける。

\paragraph{共通理論} \hyperref[theory:volcano:gases]{第\ref*{theory:volcano:gases}節}

\par 参考：\cite{vol50_gas_ref}。

\paragraph{関連問題} \hyperref[q:h28-reg-applied-28]{H28-REG-APPLIED-28}, \hyperref[q:r07-reg-applied-20]{R07-REG-APPLIED-20}

\par 難易度 4（校準初稿）、優先度未評価。

\subsection{令和7年度 問20}
\label{q:r07-reg-applied-20}

\noindent\textbf{原典：} sources/R07/17\_applied\_science.pdf、PDF物理7頁。

\textbf{公式正答：} 2（single）。

\par\textbf{検討未確定：} 理論的平均と主要の厳密な濃度基準は未提示。特にR07の100気圧・1200℃平衡の定量表と酸化還元条件の照合が残る。通常の火山ガスの代表比較として説明し、CH4の全条件不在やH2の固定割合を主張しない。

\paragraph{問題の要旨} 火山ガスの主要成分と温度・酸化還元条件を区別する。

\Needspace{4\baselineskip}\paragraph{確認した条件・記述（独自要約）}\mbox{}\par

\noindent\textbf{原典確認：} S7\_050の29原典画像で34問の全五肢・否定・数値を照合。\par

\noindent\textbf{対象：} 火山ガスの主要成分と温度・酸化還元条件を区別する。\par

\paragraph{問われる核心} 火山ガスの主要成分と温度・酸化還元条件を区別する。

\paragraph{必要知識} 水蒸気、CO2、硫黄ガス、水素ハロゲン化物等が通常の火山ガスの代表で、候補中ではCH4を選ぶ。高温マグマのガスと低温熱水・堆積物由来のガスを分ける。令和7年度は100気圧・1200℃・玄武岩質の平衡条件が原典にあり、平成28年度・令和元再試験へその条件を遡って追加しない。

\paragraph{考え方・正答を選ぶ理由（査読待ち初稿）} 水蒸気、CO2、硫黄ガス、水素ハロゲン化物等が通常の火山ガスの代表で、候補中ではCH4を選ぶ。高温マグマのガスと低温熱水・堆積物由来のガスを分ける。令和7年度は100気圧・1200℃・玄武岩質の平衡条件が原典にあり、平成28年度・令和元再試験へその条件を遡って追加しない。

\Needspace{4\baselineskip}\paragraph{選択肢の検討（査読待ち初稿）}\mbox{}\par

\noindent\textbf{1：代表的条件で適切} CO2。マグマから放出される代表的な炭素ガス。水蒸気との量比は環境に依存する。\par

\noindent\textbf{2：不適切} CH4。通常の高温マグマに由来する平均主要ガスの比較では適さない。低温・還元環境の存在を含め、厳密に常に零とする意味ではない。\par

\noindent\textbf{3：条件・解釈を留保} H2。高温ガスに含まれる還元的な水素種。量は温度・圧力だけでなく酸化還元・全体組成等にも依存する。\par

\noindent\textbf{4：代表的条件で適切} HCl。火山ガス中の代表的な水素ハロゲン化物。地下水への吸収で地表の濃度は変わり得る。\par

\noindent\textbf{5：代表的条件で適切} SO2。硫黄ガスの代表で酸化的な高温火山ガスに重要。冷却・熱水系では化学種や吸収が変わる。\par

\paragraph{再利用する考え方} 地質環境・生成機構・物質の名称を対応させ、平均値と全称、典型例と例外、当時の分類と現在の知見を分ける。

\paragraph{共通理論} \hyperref[theory:volcano:gases]{第\ref*{theory:volcano:gases}節}

\par 参考：\cite{vol50_gas_ref}。

\paragraph{関連問題} \hyperref[q:r01-ret-applied-22]{R01-RET-APPLIED-22}

\par 難易度 4（校準初稿）、優先度未評価。

\section{火山ガス・気候影響・火砕流・泥流・津波を比較する。}

\subsection{平成28年度 問35}
\label{q:h28-reg-applied-35}

\noindent\textbf{原典：} sources/H28/17\_applied\_science.pdf、PDF物理14頁。

\textbf{公式正答：} 2（single）。

\par\textbf{検討未確定：} 公式2を保持。肢1の全球冷却を大量火山灰へ帰す説明には、灰の短期影響と硫酸エアロゾルの長期影響の区別が必要。噴煙20–40kmも代表範囲。

\paragraph{問題の要旨} 火山ガス・気候影響・火砕流・泥流・津波を比較する。

\Needspace{4\baselineskip}\paragraph{確認した条件・記述（独自要約）}\mbox{}\par

\noindent\textbf{原典確認：} S7\_050の29原典画像で34問の全五肢・否定・数値を照合。\par

\noindent\textbf{対象：} 火山ガス・気候影響・火砕流・泥流・津波を比較する。\par

\paragraph{問われる核心} 火山ガス・気候影響・火砕流・泥流・津波を比較する。

\paragraph{必要知識} 火山ガスの量と有害性を分ける。通常最大成分は水蒸気なので2が誤り。火砕流はガス粒子、泥流は水粒子の混合、津波は海水の変位が中心である。1の気候説明は硫黄ガスから作る硫酸エアロゾルの長期的役割を補足し、灰だけで全球冷却を説明しない。

\paragraph{考え方・正答を選ぶ理由（査読待ち初稿）} 火山ガスの量と有害性を分ける。通常最大成分は水蒸気なので2が誤り。火砕流はガス粒子、泥流は水粒子の混合、津波は海水の変位が中心である。1の気候説明は硫黄ガスから作る硫酸エアロゾルの長期的役割を補足し、灰だけで全球冷却を説明しない。

\Needspace{4\baselineskip}\paragraph{選択肢の検討（査読待ち初稿）}\mbox{}\par

\noindent\textbf{1：条件・解釈を留保} 高い噴煙と灰による世界冷却。高い噴煙と日射遮蔽はあり得るが、長距離・長期間の全球冷却は主にSO2由来の硫酸エアロゾル。火山灰だけの説明を一般化しない。\par

\noindent\textbf{2：不適切} 大部分が有害な酸性ガス。量の大部分は通常水蒸気。少量の有害ガスの危険性と量の割合を混同している。\par

\noindent\textbf{3：代表的条件で適切} 高温ガスと粒子の火砕流。ガスと火砕粒子の混合流が高速に移動する。\par

\noindent\textbf{4：代表的条件で適切} 融雪・湖水等を取り込む火山泥流。水と火山砕屑物が混ざり、土石流状に流下する。噴火と同時以外の降雨でも発生し得る。\par

\noindent\textbf{5：代表的条件で適切} 火山島噴火や海への岩屑で津波。水塊の急激な変位により津波を起こす場合がある。必ず発生するという全称ではない。\par

\paragraph{再利用する考え方} 地質環境・生成機構・物質の名称を対応させ、平均値と全称、典型例と例外、当時の分類と現在の知見を分ける。

\paragraph{共通理論} \hyperref[theory:volcano:hazards]{第\ref*{theory:volcano:hazards}節}

\par 参考：\cite{vol50_gas_ref,vol50_climate_ref,volc_usgs_msh_flow}。

\par 難易度 4（校準初稿）、優先度未評価。

\section{パホイホイ・アア・ブロックと溶岩形態を比較する。}

\subsection{令和2年度 問24}
\label{q:r02-reg-applied-24}

\noindent\textbf{原典：} sources/R02/17\_applied\_science.pdf、PDF物理9頁。

\textbf{公式正答：} 3（single）。

\par\textbf{検討未確定：} 原典のアスペクト比1/8の分類出典・測定長さと、アアの低温・組成に関する全例の一般化は追加照合。パホイホイを単に急冷と定義しない。

\paragraph{問題の要旨} パホイホイ・アア・ブロックと溶岩形態を比較する。

\Needspace{4\baselineskip}\paragraph{確認した条件・記述（独自要約）}\mbox{}\par

\noindent\textbf{原典確認：} S7\_050の29原典画像で34問の全五肢・否定・数値を照合。\par

\noindent\textbf{対象：} パホイホイ・アア・ブロックと溶岩形態を比較する。\par

\paragraph{問われる核心} パホイホイ・アア・ブロックと溶岩形態を比較する。

\paragraph{必要知識} パホイホイとアアは同じ玄武岩質でも流動・冷却・変形の条件で異なる形になり得る。パホイホイを流紋岩質へ置換した3が誤り。滑らかな流れ、粗いクリンカー、大きな岩塊の表面を岩種と別の形態軸として整理する。

\paragraph{考え方・正答を選ぶ理由（査読待ち初稿）} パホイホイとアアは同じ玄武岩質でも流動・冷却・変形の条件で異なる形になり得る。パホイホイを流紋岩質へ置換した3が誤り。滑らかな流れ、粗いクリンカー、大きな岩塊の表面を岩種と別の形態軸として整理する。

\Needspace{4\baselineskip}\paragraph{選択肢の検討（査読待ち初稿）}\mbox{}\par

\noindent\textbf{1：条件・解釈を留保} 厚さ／長さ1/8で流れとドームを分ける。原典の形態分類の目安を保持する。普遍的な流体力学上の相転移閾値ではなく、測り方と分類文献の追加確認が必要。\par

\noindent\textbf{2：代表的条件で適切} 粘性・冷却・流速で表面が変化。組成・温度・結晶量等が流動性と表面破砕へ影響する。\par

\noindent\textbf{3：不適切} パホイホイは流紋岩質の急冷。滑らか・縄状等の表面を持つパホイホイは代表的に流動性に富む玄武岩質溶岩。岩種を交換している。\par

\noindent\textbf{4：代表的条件で適切} アアは粗いクリンカー表面。相対的に高粘性・大きい変形速度等で表面が粗い破片となる。温度だけで状態を決めない。\par

\noindent\textbf{5：代表的条件で適切} ブロックは大きな平滑破断面の岩塊。高粘性の中間質～珪長質等で、アアの小さい粗いクリンカーと区別する表面形態。\par

\paragraph{再利用する考え方} 地質環境・生成機構・物質の名称を対応させ、平均値と全称、典型例と例外、当時の分類と現在の知見を分ける。

\paragraph{共通理論} \hyperref[theory:volcano:lava]{第\ref*{theory:volcano:lava}節}

\par 参考：\cite{vol50_lava_ref,jma_magma1999}。

\par 難易度 4（校準初稿）、優先度未評価。

\section{マグマの局所生成・部分溶融・組成・発泡を比較する。}

\subsection{平成29年度 問29}
\label{q:h29-reg-applied-29}

\noindent\textbf{原典：} sources/H29/17\_applied\_science.pdf、PDF物理10頁。

\textbf{公式正答：} 3（single）。

\par\textbf{検討未確定：} 公式H29=\allowbreak{}3/R01=\allowbreak{}2。富むの比較対象が原典に明記されないため、原岩に対するNa・K等の濃集と珪長質マグマに対する組成差を区別する。発泡の絶対圧力と浮力だけによる貯留の全称も留保。

\paragraph{問題の要旨} マグマの局所生成・部分溶融・組成・発泡を比較する。

\Needspace{4\baselineskip}\paragraph{確認した条件・記述（独自要約）}\mbox{}\par

\noindent\textbf{原典確認：} S7\_050の29原典画像で34問の全五肢・否定・数値を照合。\par

\noindent\textbf{対象：} マグマの局所生成・部分溶融・組成・発泡を比較する。\par

\paragraph{問われる核心} マグマの局所生成・部分溶融・組成・発泡を比較する。

\paragraph{必要知識} 地球内部の多くは固体で、固相線を越えた条件で部分溶融する。典型的なマントルの融解液は玄武岩質であり、流紋岩等と比較すればSiO2とアルカリが少ない側なので、原典でSi・Na・Kに富むとする肢を公式は選ぶ。ただし比較を原岩マントルへ変えると不適合元素の濃集があるため、比較対象を明示して説明する。気泡の膨張と圧力の単調増加を同一視せず、マグマ貯留も中立浮力だけで全てを説明しない。

\paragraph{考え方・正答を選ぶ理由（査読待ち初稿）} 地球内部の多くは固体で、固相線を越えた条件で部分溶融する。典型的なマントルの融解液は玄武岩質であり、流紋岩等と比較すればSiO2とアルカリが少ない側なので、原典でSi・Na・Kに富むとする肢を公式は選ぶ。ただし比較を原岩マントルへ変えると不適合元素の濃集があるため、比較対象を明示して説明する。気泡の膨張と圧力の単調増加を同一視せず、マグマ貯留も中立浮力だけで全てを説明しない。

\Needspace{4\baselineskip}\paragraph{選択肢の検討（査読待ち初稿）}\mbox{}\par

\noindent\textbf{1：代表的条件で適切} 地殻・マントルは通常固体。地下全体に液状マグマが遍在するわけではない。条件が固相線を越える局所で融解が始まる。\par

\noindent\textbf{2：代表的条件で適切} 一部が融ける部分溶融。多成分の岩石では溶けやすい成分からメルトが生じ、残留固相と液相が共存する。\par

\noindent\textbf{3：条件・解釈を留保} 玄武岩質はSi・Na・Kに富む。流紋岩等と比較する通常の分類では玄武岩はSiO2・アルカリが少ない側。ただしマントル原岩に比べれば不適合元素がメルトへ濃集するため、比較対象なしの富むは曖昧。\par

\noindent\textbf{4：条件・解釈を留保} 浮力上昇とマグマだまり。通常メルトは周囲より低密度で上昇し得るが、貯留は中立浮力だけでなく応力・層境界・冷却等にも影響される。\par

\noindent\textbf{5：条件・解釈を留保} 減圧で揮発性分が分離し圧力が高まる。減圧で溶解度が下がり気泡が生じ膨張する。閉塞等で過圧を作る場合はあるが、外圧低下からガスの絶対圧力が必ず上昇するとは限らない。\par

\paragraph{再利用する考え方} 地質環境・生成機構・物質の名称を対応させ、平均値と全称、典型例と例外、当時の分類と現在の知見を分ける。

\paragraph{共通理論} \hyperref[theory:volcano:melting]{第\ref*{theory:volcano:melting}節}

\par 参考：\cite{vol50_melt_ref,jma_magma1999,vol50_degassing_ref}。

\paragraph{関連問題} \hyperref[q:r01-reg-applied-23]{R01-REG-APPLIED-23}

\par 難易度 4（校準初稿）、優先度未評価。

\subsection{令和元年度 問23}
\label{q:r01-reg-applied-23}

\noindent\textbf{原典：} sources/R01/17\_applied\_science.pdf、PDF物理8頁。

\textbf{公式正答：} 2（single）。

\par\textbf{検討未確定：} 公式H29=\allowbreak{}3/R01=\allowbreak{}2。富むの比較対象が原典に明記されないため、原岩に対するNa・K等の濃集と珪長質マグマに対する組成差を区別する。発泡の絶対圧力と浮力だけによる貯留の全称も留保。

\paragraph{問題の要旨} マグマの局所生成・部分溶融・組成・発泡を比較する。

\Needspace{4\baselineskip}\paragraph{確認した条件・記述（独自要約）}\mbox{}\par

\noindent\textbf{原典確認：} S7\_050の29原典画像で34問の全五肢・否定・数値を照合。\par

\noindent\textbf{対象：} マグマの局所生成・部分溶融・組成・発泡を比較する。\par

\paragraph{問われる核心} マグマの局所生成・部分溶融・組成・発泡を比較する。

\paragraph{必要知識} 地球内部の多くは固体で、固相線を越えた条件で部分溶融する。典型的なマントルの融解液は玄武岩質であり、流紋岩等と比較すればSiO2とアルカリが少ない側なので、原典でSi・Na・Kに富むとする肢を公式は選ぶ。ただし比較を原岩マントルへ変えると不適合元素の濃集があるため、比較対象を明示して説明する。気泡の膨張と圧力の単調増加を同一視せず、マグマ貯留も中立浮力だけで全てを説明しない。

\paragraph{考え方・正答を選ぶ理由（査読待ち初稿）} 地球内部の多くは固体で、固相線を越えた条件で部分溶融する。典型的なマントルの融解液は玄武岩質であり、流紋岩等と比較すればSiO2とアルカリが少ない側なので、原典でSi・Na・Kに富むとする肢を公式は選ぶ。ただし比較を原岩マントルへ変えると不適合元素の濃集があるため、比較対象を明示して説明する。気泡の膨張と圧力の単調増加を同一視せず、マグマ貯留も中立浮力だけで全てを説明しない。

\Needspace{4\baselineskip}\paragraph{選択肢の検討（査読待ち初稿）}\mbox{}\par

\noindent\textbf{1：条件・解釈を留保} 減圧で揮発性分が分離し圧力が高まる。減圧で溶解度が下がり気泡が生じ膨張する。閉塞等で過圧を作る場合はあるが、外圧低下からガスの絶対圧力が必ず上昇するとは限らない。\par

\noindent\textbf{2：条件・解釈を留保} 玄武岩質はSi・Na・Kに富む。流紋岩等と比較する通常の分類では玄武岩はSiO2・アルカリが少ない側。ただしマントル原岩に比べれば不適合元素がメルトへ濃集するため、比較対象なしの富むは曖昧。\par

\noindent\textbf{3：条件・解釈を留保} 浮力上昇とマグマだまり。通常メルトは周囲より低密度で上昇し得るが、貯留は中立浮力だけでなく応力・層境界・冷却等にも影響される。\par

\noindent\textbf{4：代表的条件で適切} 一部が融ける部分溶融。多成分の岩石では溶けやすい成分からメルトが生じ、残留固相と液相が共存する。\par

\noindent\textbf{5：代表的条件で適切} 地殻・マントルは通常固体。地下全体に液状マグマが遍在するわけではない。条件が固相線を越える局所で融解が始まる。\par

\paragraph{再利用する考え方} 地質環境・生成機構・物質の名称を対応させ、平均値と全称、典型例と例外、当時の分類と現在の知見を分ける。

\paragraph{共通理論} \hyperref[theory:volcano:melting]{第\ref*{theory:volcano:melting}節}

\par 参考：\cite{vol50_melt_ref,jma_magma1999,vol50_degassing_ref}。

\paragraph{関連問題} \hyperref[q:h29-reg-applied-29]{H29-REG-APPLIED-29}

\par 難易度 4（校準初稿）、優先度未評価。

\subsection{令和5年度 問25}
\label{q:r05-reg-applied-25}

\noindent\textbf{原典：} sources/R05/17\_applied\_science.pdf、PDF物理9頁。

\textbf{公式正答：} 3（single）。

\par\textbf{検討未確定：} 密度の釣合いだけで全マグマだまりを説明しない。R07原典1の融点は多成分岩石では固相線・液相線という範囲をもつ。公式R05=\allowbreak{}3/R07=\allowbreak{}2を保持して条件を補足。

\paragraph{問題の要旨} 減圧融解・含水固相線低下・部分融解を区別する。

\Needspace{4\baselineskip}\paragraph{確認した条件・記述（独自要約）}\mbox{}\par

\noindent\textbf{原典確認：} S7\_050の29原典画像で34問の全五肢・否定・数値を照合。\par

\noindent\textbf{対象：} 減圧融解・含水固相線低下・部分融解を区別する。\par

\paragraph{問われる核心} 減圧融解・含水固相線低下・部分融解を区別する。

\paragraph{必要知識} 海嶺等では上昇するマントルが減圧され、温度圧力経路が固相線を越えて融解する。低温物質の急な昇温だけを原因とした肢を選ぶ。沈み込みでは含水物質の脱水流体が固相線を下げる。純物質の一つの融点とは違い、岩石の部分溶融では固相とメルトが共存する。浮力上昇と貯留には周囲の構造や応力も関係する。

\paragraph{考え方・正答を選ぶ理由（査読待ち初稿）} 海嶺等では上昇するマントルが減圧され、温度圧力経路が固相線を越えて融解する。低温物質の急な昇温だけを原因とした肢を選ぶ。沈み込みでは含水物質の脱水流体が固相線を下げる。純物質の一つの融点とは違い、岩石の部分溶融では固相とメルトが共存する。浮力上昇と貯留には周囲の構造や応力も関係する。

\Needspace{4\baselineskip}\paragraph{選択肢の検討（査読待ち初稿）}\mbox{}\par

\noindent\textbf{1：代表的条件で適切} 昇温又は減圧で融解。温度圧力経路が固相線を越えると融解が始まる。全圧力域で同一勾配の融点を仮定しない。\par

\noindent\textbf{2：代表的条件で適切} 残留固相のある部分溶融。一部の成分が融け、他の固相が残る多成分岩石の挙動。\par

\noindent\textbf{3：不適切} 低温マントルの急な昇温が海嶺・ホットスポット成因。代表的には高温マントルの上昇に伴う減圧が融解を促す。ホットスポットの熱異常はあるが、冷たい物質の急加熱だけという説明ではない。\par

\noindent\textbf{4：代表的条件で適切} スラブ由来水による固相線低下。沈み込む含水物質の脱水で流体が供給され、上のマントルの融解が促進される。\par

\noindent\textbf{5：条件・解釈を留保} 低密度メルトの浮力と貯留。浮力は重要だが、貯留の位置は密度差だけでなく応力・構造・結晶化等に依存する。\par

\paragraph{再利用する考え方} 地質環境・生成機構・物質の名称を対応させ、平均値と全称、典型例と例外、当時の分類と現在の知見を分ける。

\paragraph{共通理論} \hyperref[theory:volcano:melting]{第\ref*{theory:volcano:melting}節}

\par 参考：\cite{vol50_melt_ref,vol50_subduction_ref}。

\paragraph{関連問題} \hyperref[q:r07-reg-applied-26]{R07-REG-APPLIED-26}

\par 難易度 4（校準初稿）、優先度未評価。

\subsection{令和7年度 問26}
\label{q:r07-reg-applied-26}

\noindent\textbf{原典：} sources/R07/17\_applied\_science.pdf、PDF物理10頁。

\textbf{公式正答：} 2（single）。

\par\textbf{検討未確定：} 密度の釣合いだけで全マグマだまりを説明しない。R07原典1の融点は多成分岩石では固相線・液相線という範囲をもつ。公式R05=\allowbreak{}3/R07=\allowbreak{}2を保持して条件を補足。

\paragraph{問題の要旨} 減圧融解・含水固相線低下・部分融解を区別する。

\Needspace{4\baselineskip}\paragraph{確認した条件・記述（独自要約）}\mbox{}\par

\noindent\textbf{原典確認：} S7\_050の29原典画像で34問の全五肢・否定・数値を照合。\par

\noindent\textbf{対象：} 減圧融解・含水固相線低下・部分融解を区別する。\par

\paragraph{問われる核心} 減圧融解・含水固相線低下・部分融解を区別する。

\paragraph{必要知識} 海嶺等では上昇するマントルが減圧され、温度圧力経路が固相線を越えて融解する。低温物質の急な昇温だけを原因とした肢を選ぶ。沈み込みでは含水物質の脱水流体が固相線を下げる。純物質の一つの融点とは違い、岩石の部分溶融では固相とメルトが共存する。浮力上昇と貯留には周囲の構造や応力も関係する。

\paragraph{考え方・正答を選ぶ理由（査読待ち初稿）} 海嶺等では上昇するマントルが減圧され、温度圧力経路が固相線を越えて融解する。低温物質の急な昇温だけを原因とした肢を選ぶ。沈み込みでは含水物質の脱水流体が固相線を下げる。純物質の一つの融点とは違い、岩石の部分溶融では固相とメルトが共存する。浮力上昇と貯留には周囲の構造や応力も関係する。

\Needspace{4\baselineskip}\paragraph{選択肢の検討（査読待ち初稿）}\mbox{}\par

\noindent\textbf{1：代表的条件で適切} 昇温又は減圧で融解。温度圧力経路が固相線を越えると融解が始まる。全圧力域で同一勾配の融点を仮定しない。\par

\noindent\textbf{2：不適切} 低温マントルの急な昇温が海嶺・ホットスポット成因。代表的には高温マントルの上昇に伴う減圧が融解を促す。ホットスポットの熱異常はあるが、冷たい物質の急加熱だけという説明ではない。\par

\noindent\textbf{3：代表的条件で適切} スラブ由来水による固相線低下。沈み込む含水物質の脱水で流体が供給され、上のマントルの融解が促進される。\par

\noindent\textbf{4：代表的条件で適切} 残留固相のある部分溶融。一部の成分が融け、他の固相が残る多成分岩石の挙動。\par

\noindent\textbf{5：条件・解釈を留保} 低密度メルトの浮力と貯留。浮力は重要だが、貯留の位置は密度差だけでなく応力・構造・結晶化等に依存する。\par

\paragraph{再利用する考え方} 地質環境・生成機構・物質の名称を対応させ、平均値と全称、典型例と例外、当時の分類と現在の知見を分ける。

\paragraph{共通理論} \hyperref[theory:volcano:melting]{第\ref*{theory:volcano:melting}節}

\par 参考：\cite{vol50_melt_ref,vol50_subduction_ref}。

\paragraph{関連問題} \hyperref[q:r05-reg-applied-25]{R05-REG-APPLIED-25}

\par 難易度 4（校準初稿）、優先度未評価。

\section{火山用語と軟堆積物の火炎状構造を区別する。}

\subsection{平成29年度 問23}
\label{q:h29-reg-applied-23}

\noindent\textbf{原典：} sources/H29/17\_applied\_science.pdf、PDF物理8頁。

\textbf{公式正答：} 4（single）。

\paragraph{問題の要旨} 火山用語と軟堆積物の火炎状構造を区別する。

\Needspace{4\baselineskip}\paragraph{確認した条件・記述（独自要約）}\mbox{}\par

\noindent\textbf{原典確認：} S7\_050の29原典画像で34問の全五肢・否定・数値を照合。\par

\noindent\textbf{対象：} 火山用語と軟堆積物の火炎状構造を区別する。\par

\paragraph{問われる核心} 火山用語と軟堆積物の火炎状構造を区別する。

\paragraph{必要知識} 火炎状構造の火炎は見た目の比喩であり、成因は未固結堆積物の変形である。火山活動の直接用語として最も弱い4を選ぶ。火砕堆積物でも変形構造ができ得るので、火山地域に絶対に存在しないという意味ではない。

\paragraph{考え方・正答を選ぶ理由（査読待ち初稿）} 火炎状構造の火炎は見た目の比喩であり、成因は未固結堆積物の変形である。火山活動の直接用語として最も弱い4を選ぶ。火砕堆積物でも変形構造ができ得るので、火山地域に絶対に存在しないという意味ではない。

\Needspace{4\baselineskip}\paragraph{選択肢の検討（査読待ち初稿）}\mbox{}\par

\noindent\textbf{1：代表的条件で適切} カルデラ。マグマの大量放出等に伴う陥没などで生じる大きな凹地で、火山活動と関係する。\par

\noindent\textbf{2：代表的条件で適切} 安山岩。島弧で代表的に見られる中間質火山岩。\par

\noindent\textbf{3：代表的条件で適切} アグルチネート。高温で放出された溶岩片等が付着・溶結してできる火砕岩。\par

\noindent\textbf{4：不適切} 火炎状構造。未固結の軽い層が上の重い堆積物の間へ押し上げられた荷重変形などで生じ、火炎という形の名は噴火を意味しない。\par

\noindent\textbf{5：代表的条件で適切} マグマ水蒸気爆発。マグマと外来水の相互作用を伴う爆発的噴火。\par

\paragraph{再利用する考え方} 地質環境・生成機構・物質の名称を対応させ、平均値と全称、典型例と例外、当時の分類と現在の知見を分ける。

\paragraph{共通理論} \hyperref[theory:volcano:terms]{第\ref*{theory:volcano:terms}節}

\par 参考：\cite{vol50_terms_ref}。

\par 難易度 3（校準初稿）、優先度未評価。

\section{噴出物の全世界体積率と火山の個数を区別する。}

\subsection{平成28年度 問29}
\label{q:h28-reg-applied-29}

\noindent\textbf{原典：} sources/H28/17\_applied\_science.pdf、PDF物理12頁。

\textbf{公式正答：} 4（single）。

\paragraph{問題の要旨} 噴出物の全世界体積率と火山の個数を区別する。

\Needspace{4\baselineskip}\paragraph{確認した条件・記述（独自要約）}\mbox{}\par

\noindent\textbf{原典確認：} S7\_050の29原典画像で34問の全五肢・否定・数値を照合。\par

\noindent\textbf{対象：} 噴出物の全世界体積率と火山の個数を区別する。\par

\paragraph{問われる核心} 噴出物の全世界体積率と火山の個数を区別する。

\paragraph{必要知識} 単位はkm3／年で、火山の個数・噴火回数・陸上被害の大きさではない。海嶺拡大が作る全世界の新しい海洋地殻に対応する4が最大。USGSの約4分の3は代表的推計で、厳密な年度別観測値や未来の一定比率ではない。

\paragraph{考え方・正答を選ぶ理由（査読待ち初稿）} 単位はkm3／年で、火山の個数・噴火回数・陸上被害の大きさではない。海嶺拡大が作る全世界の新しい海洋地殻に対応する4が最大。USGSの約4分の3は代表的推計で、厳密な年度別観測値や未来の一定比率ではない。

\Needspace{4\baselineskip}\paragraph{選択肢の検討（査読待ち初稿）}\mbox{}\par

\noindent\textbf{1：不適切} 島弧・陸弧。爆発的な陸上噴火が目立つが、全世界の体積噴出率の最大域ではない。\par

\noindent\textbf{2：不適切} 大陸地溝帯。伸張に伴う火山活動はあるが、長大な海嶺系の地殻生産に比べて小さい。\par

\noindent\textbf{3：不適切} 縁海。活動的な背弧海盆等もあるが全縁海が活動中の火山域というわけではない。\par

\noindent\textbf{4：代表的条件で適切} 中央海嶺の軸。全世界の長大な拡大軸で連続的に海洋地殻を作り、USGS概説でも年噴出物の約4分の3と見積もる。\par

\noindent\textbf{5：不適切} 海山・火山島。ハワイ等は大きな火山を作るが、個別の大型火山と全世界海嶺系の総量を比較する。\par

\paragraph{再利用する考え方} 地質環境・生成機構・物質の名称を対応させ、平均値と全称、典型例と例外、当時の分類と現在の知見を分ける。

\paragraph{共通理論} \hyperref[theory:volcano:arc]{第\ref*{theory:volcano:arc}節}

\par 参考：\cite{vol50_ridge_ref}。

\par 難易度 3（校準初稿）、優先度未評価。

\section{火砕密度流の形成と地形に沿う分布}

\subsection{令和元年度 問25}
\label{q:r01-reg-applied-25}

\noindent\textbf{原典：} sources/R01/17\_applied\_science.pdf、PDF物理9頁。

\textbf{公式正答：} 5（single）。

\par\textbf{検討未確定：} 大規模火砕流/プリニー/カルデラの一般化とプレー型の用語範囲は全類型の一次資料照合が未完了である。公式/独立5と確認済みの非単純分布を保持し部分校正に残す。

\paragraph{問題の要旨} 火砕流の発生・混合物・堆積性状と、地形を反映する分布を比較する。

\Needspace{4\baselineskip}\paragraph{確認した条件・記述（独自要約）}\mbox{}\par

\noindent\textbf{判断方向：} 最も不適切・誤りを選ぶ\par

\noindent\textbf{原典照合：} 15原典画像で18問の問題幹、全5肢、否定、専門語、化学式、候補順を確認。\par

\noindent\textbf{適用条件：} 流れと構造の典型対応は一意成因を保証しない。用語、時間間隙、保存、試料処理の条件を個別に示す。\par

\paragraph{問われる核心} 火砕流の発生・混合物・堆積性状と、地形を反映する分布を比較する。

\paragraph{必要知識} 火砕密度流の粒子と気体、ドーム崩落、噴煙柱、地形と分布、降下と流下の区別。

\paragraph{考え方・正答を選ぶ理由（査読待ち初稿）} 火砕流は地面に沿う密度流で、経路と堆積域は谷や尾根、流れの濃度に依存する。USGSには重なった扇状・舌状・ローブ状の分布及びメラピの谷埋積と希薄サージ分離の報告がある。従って単純な長円形が一般的だとする5が最も不適切である。降下火山灰は風に運ばれるため別の分布モデルを使う。大規模噴出・カルデラやプレー型という呼称は、この資料確認だけで全例を普遍化せず原典の典型説明として保持する。

\Needspace{4\baselineskip}\paragraph{選択肢の検討（査読待ち初稿）}\mbox{}\par

\noindent\textbf{1：適切（条件を含めて検討）} 一般に分級が悪く塊状となる。多粒径を混合して急速に堆積する流れの代表的特徴である。希薄サージ等の全堆積物へ塊状を強制しない。\par

\noindent\textbf{2：適切（条件を含めて検討）} 大規模火砕流は高発泡のプリニー式に伴い通常カルデラを作る。高発泡マグマの大規模噴出とカルデラ形成の典型を述べるが、全活動様式の必要十分条件としては確認していない。\par

\noindent\textbf{3：適切（条件を含めて検討）} 高温火砕物・火山ガス・取り込む大気が地面に沿って流れる。高温粒子と気体の混合物が地形に沿って流下する現象に対応する。固体だけ又は溶岩だけではない。\par

\noindent\textbf{4：適切（条件を含めて検討）} ドームの爆発型をプレー型、重力崩落型をメラピ型とする。メラピのドーム崩落起源は一次論文で確認した。プレー型という原典呼称と全類型の整理は追加確認を要する。\par

\noindent\textbf{5：誤り・最も不適切} 分布は火口を長軸の一端に持つ単純な長円が多い。谷沿いの舌状・ローブ状、複数谷の埋積などを取り得る。風で伸びる降下物の模式分布と混同する一般化である。\par

\paragraph{再利用する考え方} 観察できる構造と成因・環境の解釈を区別する。代表モデルに適用条件を添え、単一の形や測定値から唯一の環境を断定しない。

\paragraph{共通理論} \hyperref[theory:volcanology:pyroflow]{第\ref*{theory:volcanology:pyroflow}節}, \hyperref[theory:volcanology:source]{第\ref*{theory:volcanology:source}節}, \hyperref[theory:volcanology:distribution]{第\ref*{theory:volcanology:distribution}節}

\par 参考：\cite{volc_usgs_merapi,volc_usgs_msh_flow}。

\par 難易度 3（校準初稿）、優先度未評価。

